\section{The STRIVE architecture}

\subsection{Finite observations as the common language}

Many mathematical objects are infinitary when presented extensionally: a real
number has arbitrarily accurate rational approximations, a continuous function
has values at infinitely many inputs, an algebraic branch has infinitely many
compatible truncations, and a model or strategy may respond to every possible
challenge.  A finite proof does not consume the entire object.  It asks finitely
many questions and uses finitely many answers.

STRIVE represents a concept by a \emph{finite observation interface}.  For an
input parameter \(a\), a query \(q\), and a certificate \(c\), the interface
contains primitive-recursive data
\[
 B(a,q),\qquad C(a,q,c)\in\{0,1\},\qquad O(a,q,c),
\]
where \(c<B(a,q)\) is the declared finite code range, \(C=1\) means that the
certificate answers the query correctly, and \(O\) is the part of the answer
visible to the consumer.  The producer contract is a formula such as
\[
  \forall q<Q(a)\;\exists c<B(a,q)\; C(a,q,c)=1.
\]
A theorem-specific proof may establish this contract outright.  If it cannot,
the contract remains an explicit residual hypothesis.

This interface covers several familiar situations.

\begin{itemize}
\item For an exact finite object, \(q\) may be trivial and \(O=c\).
\item For a real number, \(q\) is a requested precision and \(O\) is a rational
interval or approximation.
\item For a root of a polynomial, \(c\) may contain a box and a local
verification trace, while \(O\) is a symmetric polynomial of all roots in the
box.
\item For a bounded game, \(q\) identifies a node and \(c\) is a signed
strategy code.
\item For a finite cover, \(c\) identifies a covering patch and carries a
local estimate.
\end{itemize}

The important point is that the public theorem need not expose \(c\).  The
scope of a certificate is a design choice, formalized in Section~8.

\subsection{Five layers and their proof obligations}

The framework has five layers.

\begin{description}[leftmargin=2.3cm,style=nextline]
\item[Arithmetic layer]
Natural-number coding, bounded sums/products, quotient/remainder, pairing,
order, and least search.  This is mostly inherited checked infrastructure.

\item[Strategy layer]
A compiler from each fixed bounded formula to a Boolean value, two code spaces,
two checkers, and canonical signed codes.  Mixed-radix tables are the key new
realization device.

\item[Interaction layer]
A finite query graph or clocked state machine.  It computes the set of demanded
observations, verifies compatible transitions, and replays a successful finite
interaction.  Monotone cyclic demands are first closed on a finite product of
chains.

\item[Observation layer]
Rules for erasing noncanonical certificates.  Exact observers are independent
of the certificate, quotient observers are equivalent, and metric observers
have a bounded diameter.

\item[Boundary layer]
All theorem-specific inputs are formulas in an open \(\PRAAx\) context.  A
compiled theorem has the literal shape
\[
  R_1,\ldots,R_s,\Gamma\;\vdash_{\PRAAx}\;\Phi.
\]
Supplying derivations of the residuals closes the theorem by simultaneous cut.
\end{description}

\begin{center}
\begin{tikzpicture}[
  node distance=7mm,
  box/.style={draw,rounded corners,minimum width=11cm,minimum height=9mm,
    align=center,fill=softblue},
  arr/.style={-{Latex[length=2.5mm]},thick,deepblue}]
\node[box] (arith) {Checked arithmetic: order, sums/products, division, search};
\node[box,below=of arith] (strategy) {Signed bounded-formula compiler and mixed-radix tables};
\node[box,below=of strategy] (interaction) {Demand closure, traces, viability, finite replay};
\node[box,below=of interaction] (observer) {Local, exact, quotient, or metric observation};
\node[box,below=of observer] (boundary) {Open \(\Prf\;\PRAAx\) endpoint with explicit residuals};
\draw[arr] (arith) -- (strategy);
\draw[arr] (strategy) -- (interaction);
\draw[arr] (interaction) -- (observer);
\draw[arr] (observer) -- (boundary);
\end{tikzpicture}
\end{center}

\subsection{Two kinds of compilation}

The word ``compiler'' is used in two related senses.

First, the \emph{formula compiler} is structural.  Given a fixed bounded
formula, it constructs the value and signed-certificate package.  Its
metatheoretic recursion follows the syntax tree.

Second, the \emph{contract compiler} is graph based.  Given finitely many local
lemmas and data dependencies, it determines which queries are demanded,
instantiates local results, and composes their derivations.  Its object-level
work is a clocked recurrence or a finite sequence of cuts.

The two compilers meet at a simple interface: every local finite claim can be
turned into a Boolean reflected atom or a bounded formula, and every unresolved
local producer is retained as a boundary formula.  This separation prevents a
large monolithic universal checker.  It also lets arithmetic lemmas be checked
and reused independently of any particular application.

\subsection{Why signed certificates are preferable to truth values alone}

A Boolean evaluator answers whether a finite statement is true.  A proof
compiler needs more.

\begin{itemize}
\item When the value is one, it needs enough information to introduce every
existential and disjunction and to answer every universal or conjunction
challenge.
\item When the value is zero, it needs enough information to prove an
intuitionistic negation.  A bare zero is not a proof of failure; a
counterstrategy is.
\item For diagnosis, a negative certificate identifies a least failing query or
contains a table closing every attempted witness.
\item For composition, the checker can consume an externally supplied
certificate without trusting the producer that constructed it.
\end{itemize}

The positive and negative code spaces are therefore first-class and generally
have different sizes.  Their recurrences are not implementation accidents;
they are the numerical reflection of intuitionistic introduction and
elimination behavior.
